\documentclass[english,hebrew]{article}
\usepackage[HE8,T1]{fontenc}
\usepackage[utf8]{inputenc}
\setcounter{secnumdepth}{-2}
\setlength{\parindent}{0bp}
\usepackage{xcolor}
\usepackage{calc}
\usepackage{amsmath}
\usepackage{amsthm}
\usepackage{amssymb}

\makeatletter
\@ifundefined{date}{}{\date{}}
\usepackage{graphicx}
\usepackage{geometry}
\newcommand{\limi}{\underset{n\rightarrow\infty}{\lim}}
\newcommand{\limxz}{\underset{x\rightarrow x_{0}}{\lim}}
\newcommand{\limfxz}{\underset{x\rightarrow x_{0}}{\lim}f(x)}
\newcommand{\limfxm}{\underset{x\rightarrow x_{0}^{-}}{\lim}f(x)}
\newcommand{\limfxp}{\underset{x\rightarrow x_{0}^{+}}{\lim}f(x)}
\newcommand{\limxm}{\underset{x\rightarrow x_{0}^{-}}{\lim}}
\newcommand{\limxp}{\underset{x\rightarrow x_{0}^{+}}{\lim}}
\newcommand{\sqrm}{M_{m\times m}(\mathbb{F})}
\newcommand{\seqa}{(a_{n})_{n=1}^{\infty}}
\newcommand{\seqx}{(x_{n})_{n=1}^{\infty}}
\newcommand{\funcdr}{f:D\rightarrow\mathbb{R}}
\newcommand{\der}{\limxz\frac{f(x)-f(x_{0})}{x-x_{0}}}
\usepackage[all]{pdfprivacy}

\makeatother

\usepackage{babel}
\begin{document}
\newgeometry{top=1cm, bottom=2cm, left=1cm, right=1cm}
\title{אלגברה ליניארית {\beginL 2\endL} - הרצאה {\beginL 6\endL}}
\author{נעם לב}
\date{\L{$28.10.2026$}}

\maketitle
\noindent\fbox{\begin{minipage}[t]{1\columnwidth - 2\fboxsep - 2\fboxrule}%
\begin{center}
האוניברסיטה העברית \textbar{} אלגברה ליניארית ){\beginL 2\endL}( -
\L{$80135$} \textbar{} סמסטר{\small{} }{\beginL 2027\endL}{\small{} }א'
)תשפ\textquotedblleft ז א'( \textbar{} מרצה: פרופסור יבגני סטרחוב
\textbar{} מתרגל: מר לזר גוטרמן
\par\end{center}%
\end{minipage}}
\begin{center}
הסיכום הוא לא רשמי ולא עבר אישור של הסגל, על כן אין להשתמש בו )ובין
השאר במספור שלו ובשמות הטענות שבו( כמקור רשמי. אם ישנן טעויות, הן
שלי ולא של הסגל.
\par\end{center}

\medskip{}

\tableofcontents{}

\newpage{}

\section{משפט קיילי-המילטון - \L{$\text{The Cayley-Hamilton Theorem}$}}

יהי \L{$V$} מ\textquotedblright ו מעל שדה \L{$\mathbb{F}$} כלשהו
כך ש- \L{%
\[
\dim V=n\geq1,\quad T\in\mathcal{L}(V)
\]
}%
ויהי \L{$\mathcal{B}_{V}$} בסיס של \L{$V$}.

נסמן ב- \L{$P_{ch}$} את הפולינום האופייני של \L{$T$} \L{$\left(\text{ch stands for \textbf{ch}aracteristic polynomial}\right)$}.

דהיינו,\L{%
\[
P_{ch}(x)=\det\left(I_{n}x-\underset{A}{\underbrace{M(T,\mathcal{B}_{V},\mathcal{B}_{V})}}\right)
\]
}%

כפי שהגדרנו בהרצאה {\beginL 2\endL}, \L{$P_{ch}$} הוא פולינום שאינו
תלוי בבחירת הבסיס, עם מקדמים ב-\L{$\mathbb{F}$} וממעלה לכל היותר
\L{$n$}.

לכן קיימים \L{$a_{n},a_{n-1},\dots,a_{0}\in\mathbb{F}$} כך ש-\L{%
\[
P_{ch}(x)=a_{n}x^{n}+a_{n-1}x^{n-1}+\dots+a_{1}x+a_{0}
\]
}%

בנוסף )זה לא חלק ממה שהמשפט אומר, זו עובדה(, \L{$a_{n}=1$} לכן\L{%
\[
P_{ch}(x)=x^{n}+a_{n-1}x^{n-1}+\dots+a_{1}x+a_{0}
\]
}%

אז:\L{%
\[
P_{ch}(T)=0_{\mathcal{L}(V)}
\]
}%

\noindent\fbox{\begin{minipage}[t]{1\columnwidth - 2\fboxsep - 2\fboxrule}%
בהרצאה {\beginL 2\endL} המחשנו מדוע המקדם המוביל הוא {\beginL 1\endL}. 

אנסה להראות זאת כאן שוב קצת יותר פורמאלית, על אף שיבגני התייחס לזה
כבר כעובדה.

לפי פיתוח לפלס לפי השורה הראשונה,\L{%
\[
P_{ch}(x)=\det\left(I_{n}x-\underset{A}{\underbrace{M(T,\mathcal{B}_{V},\mathcal{B}_{V})}}\right)=\begin{vmatrix}x-a_{1,1} & -a_{1,2} & \dots & -a_{1,n}\\
-a_{2,1} & x-a_{2,2} & \dots & -a_{2,n}\\
\vdots &  & \ddots & \vdots\\
-a_{n,1} & -a_{n,2} & \dots & x-a_{n,n}
\end{vmatrix}=
\]
}%
\L{%
\[
\left(x-a_{1,1}\right)\cdot\begin{vmatrix}x-a_{2,2} & \dots & -a_{2,n}\\
\vdots & \ddots & \vdots\\
-a_{n,2} & \dots & x-a_{n,n}
\end{vmatrix}+a_{1,2}\cdot\begin{vmatrix}-a_{2,1} & -a_{2,3} & \dots & -a_{2,n}\\
\vdots &  &  & \vdots\\
-a_{n,1} & -a_{n,3} & \dots & x-a_{n,n}
\end{vmatrix}+
\]
}%
\L{%
\[
\dots+\left(-1\right)^{1+n}\cdot\left(-a_{1,n}\right)\begin{vmatrix}-a_{2,1} & x-a_{2,2} & \dots & -a_{2,n-1}\\
\vdots &  &  & \vdots\\
-a_{n,1} & -a_{n,2} & \dots & -a_{n,n-1}
\end{vmatrix}
\]
}%

המחוברים כאן הם כשלעצמם פולינומים.

ניתן להוכיח באינדוקציה בעזרת פיתוח לפלס שמעלת הפולינומים למעט המחובר
הראשון היא לכל היותר \L{$n-1$} ושמעלת המחובר הראשון היא \L{$n$}.

כאשר ממשיכים לפתח בעזרת פיתוח לפלס את המחובר הראשון, היות והמקדם של
כל \L{$x$} באלכסון הוא {\beginL 1\endL}, מקבלים \L{$1\cdot x^{n}$}
כנדרש.

ניתן להוכיח גם זאת באינדוקציה.%
\end{minipage}}

\medskip{}

אז \L{$P_{ch}(T)$} הוא אופרטור על \L{$V$}, כי חזקה של אופרטור, כפל
אופרטור בסקלר וסכום אופרטורים הם אופרטורים. 

התוכן של המשפט הוא שזה בדיוק אופרטור האפס.

\subsection{\textcolor{violet}{\textquotedblright הוכחה\textquotedblright{} שגויה}}

\textcolor{violet}{יבגני רצה שנבין גם איך עלולים לשגות כאן.}

\textcolor{violet}{אפשר לנסות להוכיח בעזרת \L{%
\begin{align*}
P_{ch}\left(M(T,\mathcal{B}_{V},\mathcal{B}_{V})\right) & =\det\left(I_{n}M(T,\mathcal{B}_{V},\mathcal{B}_{V})-M(T,\mathcal{B}_{V},\mathcal{B}_{V})\right)\\
 & =\det\begin{pmatrix}0_{\mathbb{F}} & \cdots & 0_{\mathbb{F}}\\
\vdots &  & \vdots\\
0_{\mathbb{F}} & \cdots & 0_{\mathbb{F}}
\end{pmatrix}=0_{\mathbb{F}}
\end{align*}
}%
}

\textcolor{violet}{אבל זו הוכחה לא נכונה.}

\textcolor{violet}{רמז ראשוני לכך הוא שהטענה במשפט היא על שוויון של
העתקות לינאריות )או של מטריצות(, ואילו כאן באגף ימין מתקבל סקלר, איבר
בשדה.}

\textcolor{violet}{\L{$P_{ch}(A)$} מוגדר ע\textquotedblright י הצבת
\L{$A$} בפולינום האופייני של האופרטור, כלומר אחרי שהמקדמים של \L{$x^{n}+a_{n-1}x^{n-1}+\dots+a_{1}x+a_{0}$}
כבר חושבו.}

\textcolor{violet}{כאן עושים משהו אחר: מחשבים את הדטרמיננטה של ביטוי
כלשהו )הביטוי \L{$I_{n}M(T,\mathcal{B}_{V},\mathcal{B}_{V})-M(T,\mathcal{B}_{V},\mathcal{B}_{V})$}(,
וזה כמובן יוצא סקלר - לעומת חישוב של הפולינום האופייני ולאחר מכן הצבה
של המטריצה בו. }

\subsection{הוכחת המשפט}
\begin{proof}
1. נסמן \L{$M(T,\mathcal{B}_{V},\mathcal{B}_{V})=A$}. מספיק להוכיח
כי \L{%
\[
P_{ch}(A)=A^{n}+a_{n-1}A^{n-1}+\dots+a_{1}A+a_{0}I_{n}=0
\]
}%

כיוון שיש איזומורפיזם בין מרחב האופרטורים למרחב המטריצות, ולכן שוויון
אלגברי במרחב המטריצות נותן את אותו שוויון במרחב האופרטורים, במיוחד
כש-\L{$A$} היא המטריצה של \L{$T$} בבסיס הנתון.

\textcolor{teal}{\emph{התמקדתי בזה בלמה בסוף סיכום הרצאה \L{$4$}.}}

\textcolor{teal}{\emph{עבור \L{$\Phi(S)=M(S,\mathcal{B}_{V},\mathcal{B}_{V})$}
ראינו ש- \L{$\Phi(q(T))=q(A)$} לכל פולינום \L{$q$} )ובפרט עם \L{$q=P_{ch}$}(,
ולכן \L{$q(T)=0\iff q(A)=0$}.}}

2. נזכיר את הנוסחה מהרצאה \L{$5$}: אם \L{%
\[
Q=\left(q_{i,j}\right)_{i,j=1}^{n}
\]
}%
הפיכה, אז \L{%
\[
Q^{-1}=\left(d_{i,j}\right)_{i,j=1}^{n}\cdot\frac{1}{\det(Q)}
\]
}%
כאשר לכל \L{$1\leq i,j\leq n$} \L{%
\[
d_{i,j}=(-1)^{i+j}\cdot\det\left(\begin{array}{cc|c|cc}
q_{1,1} & \cdots & q_{1,i} & \cdots & q_{1,n}\\
\vdots &  & \vdots &  & \vdots\\
\hline q_{j,1} & \cdots & q_{j,i} & \cdots & q_{j,n}\\
\hline \vdots &  & \vdots &  & \vdots\\
q_{n,1} & \cdots & q_{n,i} & \cdots & q_{n,n}
\end{array}\right)=(-1)^{i+j}\cdot\det\left(Q_{j,i}\right)
\]
}%
)מוחקים את השורה ה-\L{$j$} ואת העמודה ה-\L{$i$}: הקווים מסמנים את
השורה ה-\L{$j$} ואת העמודה ה-\L{$i$}, שאותן מוחקים; \L{$Q_{j,i}$}
הוא המינור המתאים(.

נגדיר \L{%
\[
Q\overset{\text{def}}{=}I_{n}x-A
\]
}%
\L{%
\[
\text{adj}(Q)\overset{\text{def}}{=}\left(d_{i,j}\right)_{i,j=1}^{n}
\]
}%
כלומר \L{%
\[
Q^{-1}=\text{adj}(Q)\cdot\frac{1}{\det(Q)}
\]
}%
ולכן \L{%
\[
\text{adj}(Q)\cdot Q=I_{n}\cdot\text{adj}(Q)\cdot Q=\det(Q)\cdot\frac{1}{\det(Q)}\cdot\text{adj}(Q)\cdot Q=\det(Q)\cdot Q^{-1}\cdot Q=\det(Q)\cdot I_{n}
\]
}%
\L{%
\[
Q\cdot\text{adj}(Q)=Q\cdot\text{adj}(Q)\cdot I_{n}=Q\cdot\text{adj}(Q)\cdot\frac{1}{\det(Q)}\cdot\det(Q)=Q\cdot Q^{-1}\cdot\det(Q)=I_{n}\cdot\det(Q)=\det(Q)\cdot I_{n}
\]
}%

לפעמים קוראים ל- \L{$\text{adj}(Q)$} ה- \L{$\text{\textbf{adjoint}}$}
של \L{$Q$}, אבל יהיה לנו גם \L{$\text{adjoint}$} של העתקה לינארית
במובן אחר, אז כאן זו פשוט מטריצה שמוגדרת לצורך ההוכחה.

בפרט, \L{$\text{adj}(Q)$} ו-\L{$Q$} מתחלפות, וזה לא היה ברור מראש.

כאן הוכחנו זאת בעזרת שימוש בזה ש-\L{$Q$} הפיכה, אך חשוב לציין שיש
הוכחה שלא משתמשת בכך ודומה מאוד להוכחה של המשפט בסוף הרצאה {\beginL 5\endL}.

לכן השוויונות שהוכחנו עכשיו\L{%
\[
Q\cdot\text{adj}(Q)=\det(Q)\cdot I_{n}=\text{adj}(Q)\cdot Q
\]
}%
נכונים לכל מטריצה ריבועית \L{$Q$}, גם אם היא לא הפיכה.

זה רלוונטי כי \L{$Q=I_{n}x-A$} אינה הפיכה בדיוק כאשר \L{$x$} ערך
עצמי של \L{$A$}, ובשלב 3 אנחנו צריכים את השוויון לכל \L{$x$}.\\

\textcolor{teal}{סקירה של ההוכחה בדרך האחרת: }

\textcolor{teal}{השוויון הראשון הוא בדיוק מה שהוכחנו בשלבים 3 ו-4
בהוכחה מהרצאה \L{$5$}, אם לא מחלקים בדטרמיננטה: האיבר במקום \L{$(i,i)$}
של \L{$Q\cdot\text{adj}(Q)$} הוא פיתוח של \L{$\det(Q)$} לפי השורה
ה-\L{$i$}, והאיבר במקום \L{$(i,j)$} עבור \L{$i\neq j$} הוא פיתוח
של דטרמיננטה עם שתי שורות זהות, ולכן \L{$0$}.}

\textcolor{teal}{ההנחה \L{$\det(Q)\neq0$} שימשה שם רק לחלוקה בסוף.}

\textcolor{teal}{השוויון השני נובע מהראשון דרך שחלוף: מההגדרה \L{$\text{adj}(Q^{T})=\text{adj}(Q)^{T}$},
כי המינור של \L{$Q^{T}$} במקום \L{$(j,i)$} הוא השחלוף של המינור
של \L{$Q$} במקום \L{$(i,j)$}, ודטרמיננטה לא משתנה בשחלוף )לינארית
1, הרצאה \L{$29$}(.}

\textcolor{teal}{לכן מהשוויון הראשון עבור \L{$Q^{T}$}: \L{%
\[
Q^{T}\cdot\text{adj}(Q)^{T}=\det(Q^{T})\cdot I_{n}=\det(Q)\cdot I_{n}
\]
}%
ובשחלוף שני האגפים )\L{$(XY)^{T}=Y^{T}X^{T}$}( מקבלים \L{$\text{adj}(Q)\cdot Q=\det(Q)\cdot I_{n}$}.}

3. נציב \L{$Q=I_{n}x-A$}: \L{%
\[
(I_{n}x-A)\cdot\text{adj}(I_{n}x-A)=\det(I_{n}x-A)\cdot I_{n}
\]
}%

באגף ימין יש למעשה בדיוק את \L{$P_{ch}(x)\cdot I_{n}$}, כלומר בדיוק
המקדמים שאנחנו רוצים.

\textcolor{violet}{עוד \textquotedblright הוכחה\textquotedblright{}
שגויה: להציב \L{$x=A$} ולקבל \L{%
\[
0\cdot\text{adj}(I_{n}x-A)=P_{ch}(A)
\]
}%
זה נראה בדיוק כמו מה שרצינו, אבל זה לא נכון.}

\textcolor{violet}{מותר להציב מטריצה \L{$A$} בזהות פולינומית כשהמקדמים
סקלרים )הרצאה \L{$5$}(, או כשהם מטריצות המתחלפות עם \L{$A$}.}

\textcolor{violet}{בשלב הזה לא ידוע לנו דבר על המקדמים של \L{$\text{adj}(I_{n}x-A)$},
ולכן ההצבה אינה מוצדקת.}

\textcolor{violet}{)בדיעבד, משלב 6 נובע שהם כן מתחלפים עם \L{$A$},
אבל אז ההוכחה כבר הושלמה בלעדיה(.}

\textcolor{violet}{יבגני המחיש בדוגמה מה קורה כשהמקדמים לא מתחלפים:
השוויון \L{%
\[
(xI_{n}-A)(xI_{n}+B)=x^{2}I_{n}+(B-A)x-AB
\]
}%
נכון לכל סקלר \L{$x$}, אבל בהצבת \L{$x=A$} אגף שמאל מתאפס ואגף ימין
הוא \L{$BA-AB$}, שאינו \L{$0$}, למשל עבור \L{%
\[
A=\begin{pmatrix}0 & 1\\
1 & 0
\end{pmatrix},\quad B=\begin{pmatrix}1 & 0\\
0 & -1
\end{pmatrix}
\]
}%
}

{\beginL 4\endL}. כל איבר של \L{$\text{adj}(I_{n}x-A)$} הוא, עד כדי
סימן, דטרמיננטה של מטריצת המינור שהיא בגודל \L{$(n-1)\times(n-1)$}
ואיברי המינור הם פולינומים במשתנה \L{$x$} ממעלה לכל היותר \L{$1$},
ולכן הדטרמיננטה של המינור הוא פולינום במשתנה \L{$x$} ממעלה לכל היותר
\L{$n-1$}.

דהיינו, כל איבר של \L{$\text{adj}(I_{n}x-A)$} הוא פולינום במשתנה
\L{$x$} ממעלה לכל היותר \L{$n-1$}.

נפרק את \L{$\text{adj}(I_{n}x-A)$} לפי החזקות שלו: \L{%
\[
\text{adj}(I_{n}x-A)=B_{n-1}x^{n-1}+B_{n-2}x^{n-2}+\dots+B_{1}x+B_{0}
\]
}%
ולפי 3\L{%
\[
(I_{n}x-A)\left(B_{n-1}x^{n-1}+B_{n-2}x^{n-2}+\dots+B_{1}x+B_{0}\right)=x^{n}I_{n}+x^{n-1}a_{n-1}I_{n}+\dots+xa_{1}I_{n}+a_{0}I_{n}
\]
}%

נשים לב שמתקיים: \L{%
\begin{align*}
(I_{n}x-A)\left(B_{n-1}x^{n-1}+\dots+B_{1}x+B_{0}\right) & =B_{n-1}x^{n}+B_{n-2}x^{n-1}+\dots+B_{1}x^{2}+B_{0}x\\
 & \quad-AB_{n-1}x^{n-1}-AB_{n-2}x^{n-2}-\dots-AB_{1}x-AB_{0}
\end{align*}
}%
\L{%
\[
=B_{n-1}x^{n}+\left(B_{n-2}-AB_{n-1}\right)x^{n-1}+\dots+\left(B_{0}-AB_{1}\right)x-AB_{0}
\]
}%

{\beginL 5\endL}. נשווה את המקדמים לפני כל חזקה של \L{$x$} בשני האגפים:
\L{%
\begin{align*}
x^{n}: & \quad B_{n-1}=I_{n}\\
x^{n-1}: & \quad B_{n-2}-AB_{n-1}=a_{n-1}I_{n}\\
x^{n-2}: & \quad B_{n-3}-AB_{n-2}=a_{n-2}I_{n}\\
 & \quad\vdots\\
x^{2}: & \quad B_{1}-AB_{2}=a_{2}I_{n}\\
x: & \quad B_{0}-AB_{1}=a_{1}I_{n}\\
1: & \quad-AB_{0}=a_{0}I_{n}
\end{align*}
}%
מותר להשוות מקדמים: השוויון ב-{\beginL 4\endL} מתקיים לכל \L{$x\in\mathbb{F}$},
ובכל איבר \L{$(i,j)$} שני האגפים הם פולינומים ב-\L{$x$}.

מעל \L{$\mathbb{R}$} או \L{$\mathbb{C}$}, פולינום שמתאפס באינסוף
נקודות הוא פולינום האפס, ולכן המקדמים שווים איבר-איבר.

זו מערכת של \L{$n+1$} משוואות, כל אחת היא קשר בין מטריצות. 

יש כאן \textquotedblright קונספירציה\textquotedblright : המקדמים של
הפולינום האופייני תלויים במטריצה \L{$A$}, ובזכות זה מתקבל \L{$0$}.

{\beginL 6\endL}. לכן \L{%
\begin{align*}
P_{ch}(x)I_{n} & =x^{n}I_{n}+x^{n-1}(a_{n-1}I_{n})+x^{n-2}(a_{n-2}I_{n})\\
 & \quad+\dots+x^{2}(a_{2}I_{n})+x(a_{1}I_{n})+a_{0}I_{n}
\end{align*}
}%

נכפול כל משוואה מ-5 ב-\L{$A^{k}$} המתאים משמאל ונסכום: \L{%
\begin{align*}
P_{ch}(A) & =A^{n}B_{n-1}+A^{n-1}(B_{n-2}-AB_{n-1})+A^{n-2}(B_{n-3}-AB_{n-2})+\dots\\
 & \quad+A^{2}(B_{1}-AB_{2})+A(B_{0}-AB_{1})-AB_{0}\\
 & =A^{n}B_{n-1}+\left(A^{n-1}B_{n-2}-A^{n}B_{n-1}\right)+\left(A^{n-2}B_{n-3}-A^{n-1}B_{n-2}\right)+\dots\\
 & \quad+\left(A^{2}B_{1}-A^{3}B_{2}\right)+\left(AB_{0}-A^{2}B_{1}\right)-AB_{0}\\
 & =A^{n}B_{n-1}-A^{n}B_{n-1}+A^{n-1}B_{n-2}-A^{n-1}B_{n-2}+\dots+AB_{0}-AB_{0}\\
 & =0
\end{align*}
}%

זה נראה כמו ההצבה \L{$x=A$} מ-4, אבל יש הבדל: ב-4 הצבנו \L{$A$}
בזהות אחת שבה \L{$x$} מופיע גם מימין למקדמים המטריציוניים, וזה דרש
להעביר את \L{$A$} מעבר ל-\L{$B_{k}$}. 

כאן, לעומת זאת, כל משוואה מ-5 היא שוויון בין מטריצות קבועות )בלי \L{$x$}(,
ואפשר לכפול שוויון כזה ב-\L{$A^{k}$} משמאל.

כלומר, אין כאן בעיה כי לא שינינו את סדר הכפל בשום מקום; מההתחלה התייחסנו
לכפל כלא קומוטטיבי.

יבגני אמר שההוכחה לא תינתן במבחן כי היא טכנית מדי: מספיק לטעות קצת
באינדקס וכל ההוכחה נופלת.
\end{proof}

\section{מסקנה - הפולינום המינימלי מחלק את הפולינום האופייני}

יהי \L{$V$} מ\textquotedblright ו מעל שדה \L{$\mathbb{F}$} כלשהו
כך ש- \L{%
\[
\dim V=n\geq1,\quad T\in\mathcal{L}(V)
\]
}%
ויהי \L{$q_{\min}(x)$} הפולינום המינימלי של \L{$T$}. אז:

)א( \L{$q_{\min}$} מחלק את הפולינום האופייני של \L{$T$}.

)ב( \L{%
\[
1\leq\deg(q_{\min})\leq n
\]
}%

\begin{proof}
)א( לפי משפט קיילי-המילטון \L{$P_{ch}(T)=0$}, ולפי המשפט מהרצאה \L{$5$}
הפולינום המינימלי מחלק כל פולינום שמאפס את \L{$T$}.

)ב( \L{$q_{\min}$} מחלק את \L{$P_{ch}$}, ו-\L{$\deg(P_{ch})=n$},
ולכן \L{$\deg(q_{\min})\leq n$}.

בנוסף, \L{$q_{\min}$} לא יכול להיות קבוע, כי קבוע שונה מאפס כפול
\L{$I$} אינו אופרטור האפס:

דהיינו, בגלל שהפולינום המינימלי מחלק את הפולינום האופייני, קיים פולינום
\L{$S$} כך ש{\beginL -\endL} \L{$P_{ch}=S\cdot q_{\min}$}

וגם הפולינום האופייני הוא לא פולינום האפס כי המקדם המוביל שלו הוא
{\beginL 1\endL}. 

לכן \L{$S\neq0$} ומכאן\L{%
\[
n=\deg(S)+\deg(q_{\min})\geq\deg(q_{\min})
\]
}%
אם \L{$\deg(q_{\min})=0$} אז היות והפולינום המינימלי הוא פולינום
מתוקן נקבל \L{$q_{\min}=1$}.

אבל \L{$\dim V\geq1$} לכן \L{$q_{\min}(T)=I\neq0$} וזו סתירה.
\end{proof}
החסם \L{$\deg(q_{\min})\leq n$} לא ברור מאליו, והוא נובע רק מקיילי-המילטון. 

בהוכחת הקיום בהרצאה \L{$4$} קיבלנו רק \L{$\deg(q_{\min})\leq n^{2}$},
מהרשימה \L{$(I,T,\dots,T^{n^{2}})$}.

\section{משפט - הפולינום המינימלי וייצוג ע\textquotedblright י מטריצה משולשית
עליונה}

יהי \L{$V$} מ\textquotedblright ו מעל שדה \L{$\mathbb{F}$} כלשהו
כך ש- \L{%
\[
\dim V=n\geq1,\quad T\in\mathcal{L}(V)
\]
}%
ויהי \L{$q_{\min}$} הפולינום המינימלי של \L{$T$}. אז התנאים הבאים
שקולים:

\L{$(\text{I})$} קיים בסיס \L{$\mathcal{B}_{V}$} במרחב הווקטורי
\L{$V$} כך ש- \L{%
\[
M(T,\mathcal{B}_{V},\mathcal{B}_{V})=\begin{pmatrix}a_{1,1} & \cdots & a_{1,n}\\
0_{\mathbb{F}} & \ddots & \vdots\\
0_{\mathbb{F}} & \cdots0_{\mathbb{F}} & a_{n,n}
\end{pmatrix}
\]
}%

\L{$(\text{II})$} ניתן לפרק את הפולינום המינימלי לגורמים \L{%
\[
q_{\min}(z)=(z-\lambda_{1})\cdots(z-\lambda_{m})
\]
}%
כאשר \L{%
\[
\lambda_{1},\dots,\lambda_{m}\in\mathbb{F}
\]
}%

\L{$\lambda_{1},\dots,\lambda_{m}$} לא חייבים להיות שונים זה מזה;
ייתכן שחלקם, או אפילו כולם, שווים. \\

\textcolor{teal}{זה המשפט המשמעותי השני על פולינומים מינימליים.}

\textcolor{teal}{מבחינה שימושית הכיוון \L{$(\text{II})\Rightarrow(\text{I})$}
חשוב יותר: כדי לדעת אם אפשר להציג את \L{$T$} ע\textquotedblright י
מטריצה משולשית עליונה, בודקים את הפולינום המינימלי.}

\textcolor{teal}{מעל \L{$\mathbb{C}$} אנחנו יודעים מראש שזה תמיד
אפשרי, מעל \L{$\mathbb{R}$} לא.}

\textcolor{teal}{לפי המשפט מהרצאה \L{$5$}, \L{$\lambda_{1},\dots,\lambda_{m}$}
הם בדיוק הערכים העצמיים של \L{$T$} )עם חזרות(, כלומר \L{$(\text{II})$}
אומר ש-\L{$q_{\min}$} מתפרק לגורמים לינאריים מעל \L{$\mathbb{F}$}.}

\textcolor{teal}{מעל \L{$\mathbb{C}$} זה תמיד מתקיים לפי המשפט היסודי
של האלגברה )הרצאה \L{$2$}(, ולכן המשפט מהרצאה \L{$3$} )מעל \L{$\mathbb{C}$}
כל אופרטור ניתן לייצוג ע\textquotedblright י מטריצה משולשית עליונה(
הוא מקרה פרטי של המשפט הזה. }

\textcolor{teal}{מעל \L{$\mathbb{R}$}, הסיבוב \L{$T(x,y)=(-y,x)$}
מהרצאה \L{$4$} הוא דוגמה נגדית: אותו חישוב נותן \L{$q_{\min}(z)=z^{2}+1$},
שאין לו שורשים ממשיים.}

\subsection{הערה - שורשי \L{$q_{\min}$} הם איברי אלכסון}

מתקיים \L{%
\[
\lambda_{1},\dots,\lambda_{m}\in\left\{ a_{1,1},\dots,a_{n,n}\right\} 
\]
}%

כאן \L{$a_{1,1},\dots,a_{n,n}$} הם איברי האלכסון של המטריצה מ-\L{$(\text{I})$}.

למעשה זה שוויון של קבוצות כי כבר יש לנו את הכיוון השני: לפי ההערה
בהרצאה \L{$4$}, כל \L{$a_{i,i}$} הוא ערך עצמי של \L{$T$}, ולכן
לפי הרצאה \L{$5$} הוא שורש של \L{$q_{\min}$}.

\L{$m\leq n$}, ולא בהכרח \L{$m=n$}.

למשל, אם \L{$P_{ch}(z)=(z-2)^{3}(z-1)$} )כלומר \L{$n=4$}(, האפשרויות
ל-\L{$q_{\min}$} הן \L{%
\[
(z-2)^{3}(z-1),\quad(z-2)^{2}(z-1),\quad(z-2)(z-1)
\]
}%
כי כל ערך עצמי הוא שורש של הפולינום המינימלי, ולכן \L{$z-2$} ו-\L{$z-1$}
חייבים להופיע.

\subsection{הוכחת המשפט}
\begin{proof}
\L{$(\text{I})$} \L{$\Leftarrow$} \L{$(\text{II})$}

1. \L{%
\begin{align*}
\det\left(zI_{n}-M(T,\mathcal{B}_{V},\mathcal{B}_{V})\right) & =(z-a_{1,1})(z-a_{2,2})\cdots(z-a_{n,n})\\
 & =P_{ch}(z)
\end{align*}
}%

מפתחים לפי העמודה הראשונה, שבה רק האיבר הראשון שונה מאפס, וממשיכים
באינדוקציה.

זה החישוב מההערה \textquotedblright איברי האלכסון של מטריצה משולשית
עליונה הם ערכים עצמיים\textquotedblright{} בהרצאה \L{$4$}.

2. \L{$q_{\min}(z)$} מחלק את \L{$P_{ch}(z)$} לכן\L{%
\[
q_{\min}(z)=(z-\lambda_{1})\cdots(z-\lambda_{m})
\]
}%

יבגני התייחס לזה כמובן מאליו שמחלק של מכפלת גורמים לינאריים הוא מכפלה
של חלק מהם אבל ניתן להוכיח זאת באינדוקציה על הדרגה: פולינום מתוקן
שמחלק מכפלה של גורמים לינאריים הוא מכפלה של חלק מהם.

אם \L{$\deg(q_{\min})=0$} אז \L{$q_{\min}=1$} )מתוקן(, וזה מתקיים
באופן טריוויאלי.

אחרת, היות והפולינום המינימלי מחלק אותו נוכל לכתוב\L{%
\[
q_{\min}(z)S(z)=(z-a_{1,1})\cdot_{\dots}(z-a_{n,n})
\]
}%

ונציב \L{$z=a_{1,1}$}: אגף שמאל מתאפס, לכן \L{$q_{\min}(a_{1,1})=0_{\mathbb{F}}$}
או \L{$S(a_{1,1})=0_{\mathbb{F}}$}.

אם \L{$q_{\min}(a_{1,1})=0_{\mathbb{F}}$}, אז \L{$q_{\min}(z)=(z-a_{1,1})q_{1}(z)$}
כאשר \L{$q_{1}$} הוא פולינום מתוקן ממעלה \L{$\deg(q_{\min})-1$}. 

נציב בשוויון ונצמצם את \L{$(z-a_{1,1})$}: \L{%
\[
q_{1}(z)S(z)=(z-a_{2,2})\cdot_{\dots}(z-a_{n,n})
\]
}%
לפי הנחת האינדוקציה \L{$q_{1}$} הוא מכפלה של חלק מהגורמים, ולכן גם
\L{$q_{\min}=(z-a_{1,1})q_{1}$}.

אחרת \L{$S(a_{1,1})=0_{\mathbb{F}}$}, לכן \L{$S(z)=(z-a_{1,1})S_{1}(z)$}.
נצמצם: \L{%
\[
q_{\min}(z)S_{1}(z)=(z-a_{2,2})\cdot_{\dots}(z-a_{n,n})
\]
}%
כלומר \L{$q_{\min}$} מחלק מכפלה של \L{$n-1$} גורמים.

נחזור על אותו טיעון עם \L{$a_{2,2}$}, ואם שוב \L{$S_{1}$} הוא שמתאפס
נמשיך ל-\L{$a_{3,3}$} וכן הלאה.

זה חייב להיעצר במקום מסוים, כי אחרת אחרי \L{$n$} צמצומים נקבל \L{$q_{\min}(z)S_{n}(z)=1$},
ולכן \L{$\deg(q_{\min})=0$}, בסתירה להנחה.

הצמצום מותר כי אם \L{$(z-c)f=(z-c)g$} אז \L{$(z-c)(f-g)=0$}, ואם
\L{$f\neq g$} המעלה של אגף שמאל היא לפחות \L{$1$}, בסתירה.

\L{$(\text{II})$} \L{$\Leftarrow$} \L{$(\text{I})$}

באינדוקציה על \L{$m$}.

1. נניח כי \L{$m=1$}, אז \L{%
\[
q_{\min}(z)=z-\lambda_{1}
\]
}%
\L{%
\[
q_{\min}(T)=T-\lambda_{1}I=0
\]
}%
\L{%
\[
\therefore\quad T=\lambda_{1}I
\]
}%
\L{%
\[
\therefore\quad M(T,\mathcal{B}_{V},\mathcal{B}_{V})=\begin{pmatrix}\lambda_{1} &  & 0\\
 & \ddots\\
0 &  & \lambda_{1}
\end{pmatrix}
\]
}%
בכל בסיס ב-\L{$V$}, וזו מטריצה אלכסונית שזה מקרה פרטי של משולשית
עליונה.

2. יהי \L{$m\geq2$}. נניח שאם \L{%
\[
\widetilde{q}_{\min}(z)=(z-\alpha_{1})\cdots(z-\alpha_{k})
\]
}%
)כאשר \L{$1\leq k\leq m-1$}( הוא הפולינום המינימלי של \L{$\widetilde{T}\in\mathcal{L}(\widetilde{V})$},
אז קיים בסיס \L{$\mathcal{B}_{\widetilde{V}}$} ב-\L{$\widetilde{V}$}
כך ש- \L{$M(\widetilde{T},\mathcal{B}_{\widetilde{V}},\mathcal{B}_{\widetilde{V}})$}
משולשית עליונה.

זו הנחת אינדוקציה מאוד חזקה: היא על כל מרחב וקטורי \L{$\widetilde{V}$}
ועל כל אופרטור \L{$\widetilde{T}$} שהפולינום המינימלי שלו הוא מכפלה
של פחות מ-\L{$m$} גורמים לינאריים. 

ממנה נגיע ל-\L{$V$}, \L{$T$} ו-\L{$q_{\min}$} שלנו.

3. נגדיר \L{%
\[
U=\text{Range}(T-\lambda_{m}I)
\]
}%
\L{$U$} - תת-מרחב שמור של \L{$T$}:

אם \L{$u\in U$}, אז \L{%
\[
T(u)=\underset{\in\text{Range}(T-\lambda_{m}I)=U}{\underbrace{(T-\lambda_{m}I)(u)}}+\underset{\in U}{\underbrace{\lambda_{m}u}}
\]
}%
ולכן \L{$T(u)\in U$}.

{\beginL 4\endL}. יהי \L{$u\in U$}, כלומר \L{$u=(T-\lambda_{m}I)\left(v_{0}\right)$}
עבור \L{$v_{0}\in V$} כלשהו. אז

\L{%
\begin{align*}
(T-\lambda_{1}I)\cdot_{\dots}(T-\lambda_{m-1}I)\left(u\right) & =\underset{q_{\min}(T)=0}{\underbrace{(T-\lambda_{1}I)\cdot_{\dots}(T-\lambda_{m-1}I)\cdot(T-\lambda_{m}I)}}\left(v_{0}\right)\\
 & =0
\end{align*}
}%
נגדיר \L{%
\[
p(z)=(z-\lambda_{1})\cdots(z-\lambda_{m-1})
\]
}%
ואז \L{%
\[
p(T|_{U})=0
\]
}%

אם \L{$\widetilde{q}_{\min}$} הוא הפולינום המינימלי של \L{$T|_{U}$},
אז \L{$\widetilde{q}_{\min}$} מחלק את \L{$p(z)$} ולכן

\L{%
\[
\widetilde{q}_{\min}(z)=(z-\alpha_{1})\cdots(z-\alpha_{k})
\]
}%
\L{%
\[
1\leq k\leq m-1
\]
}%

5. נשתמש ב-2 עבור \L{%
\[
\widetilde{V}=U=\text{Range}(T-\lambda_{m}I)
\]
}%
\L{%
\[
\widetilde{T}=T|_{U}
\]
}%
ועבור \L{$\widetilde{q}_{\min}$}.

לכן קיים בסיס \L{$\mathcal{B}_{U}$} ב-\L{$U$} כך ש-\L{$M(T|_{U},\mathcal{B}_{U},\mathcal{B}_{U})$}
משולשית עליונה. \L{%
\[
\mathcal{B}_{U}=(u_{1},\dots,u_{p})
\]
}%
\L{%
\[
T(u_{j})\in\text{Span}(u_{1},\dots,u_{j}),\quad1\leq j\leq p
\]
}%

6. נשלים את \L{$\mathcal{B}_{U}$} לבסיס \L{$\mathcal{B}_{V}$}: \L{%
\[
\mathcal{B}_{V}=(u_{1},\dots,u_{p},v_{1},\dots,v_{s})
\]
}%
נרצה להוכיח כי \L{$M(T,\mathcal{B}_{V},\mathcal{B}_{V})$} משולשית
עליונה. מספיק להוכיח כי \L{%
\[
T(v_{\ell})\in\text{Span}(u_{1},\dots,u_{p},v_{1},\dots,v_{\ell})
\]
}%

ואכן \L{%
\begin{align*}
T(v_{\ell}) & =\underset{\in U=\text{Span}(u_{1},\dots,u_{p})}{\underbrace{(T-\lambda_{m}I)(v_{\ell})}}+\lambda_{m}v_{\ell}\in\text{Span}(u_{1},\dots,u_{p},v_{\ell})\\
 & \subseteq\text{Span}(u_{1},\dots,u_{p},v_{1},\dots,v_{\ell})
\end{align*}
}%

יחד עם 5, כל וקטור ב-\L{$\mathcal{B}_{V}$} נשלח לפרוש של עצמו ושל
קודמיו, ולכן לפי הלמה מהרצאה 3 המטריצה משולשית עליונה.
\end{proof}

\end{document}
